
\documentclass{article}
\usepackage{colortbl}
\usepackage{makecell}
\usepackage{multirow}
\usepackage{supertabular}

\begin{document}

\newcounter{utterance}

\centering \large Interaction Transcript for game `wordle', experiment `high\_frequency\_words\_no\_clue\_no\_critic', episode 9 with reflexion{-}Qwen3.5{-}27B.
\vspace{24pt}

{ \footnotesize  \setcounter{utterance}{1}
\setlength{\tabcolsep}{0pt}
\begin{supertabular}{c@{$\;$}|p{.15\linewidth}@{}p{.15\linewidth}p{.15\linewidth}p{.15\linewidth}p{.15\linewidth}p{.15\linewidth}}
    \# & \multicolumn{2}{c}{Player} && \multicolumn{2}{c}{Game Master} \\
    \hline

    \theutterance \stepcounter{utterance}  
    & & & \multicolumn{4}{p{0.6\linewidth}}{
        \cellcolor[rgb]{0.9,0.9,0.9}{
            \makecell[{{p{\linewidth}}}]{
                \texttt{\tiny{[P1$\langle$GM]}}
                \texttt{You have attempted this same task before and did not succeed. These are the plans you wrote after those attempts.} \\
\texttt{Trial \#1: New plan: The previous attempt failed because I entered a repetitive loop in my explanation, repeatedly listing the same excluded words ("Loves", "Lived", "Leave") without generating a new valid guess, which caused the response to be truncated and rejected. I failed to strictly adhere to the "concise explanation" rule and the requirement to provide a valid 5{-}letter word. Next time, when stuck on a pattern like "l\_ve\_", I will immediately stop listing invalid candidates and directly propose a single, valid English word that fits the constraints (e.g., "lives" if 's' wasn't red, or "loved" if 'o'/'d' weren't red, or a new candidate like "lurve" if valid, or "laved" if 'a' wasn't red). I will ensure the explanation is brief, explicitly states the logic for the *new* guess, and ends with a valid 5{-}letter word in the "guess:" tag, avoiding any repetition or rambling.} \\
\\ 
\texttt{Trial \#2: New plan: The previous attempt failed because I entered a repetitive loop in my explanation, listing numerous invalid or non{-}existent words (e.g., "Vilem", "Vilen") instead of generating a valid 5{-}letter guess, which caused the response to be truncated and rejected. I failed to strictly adhere to the "concise explanation" rule and the requirement to provide a valid English word. Next time, when stuck on a pattern like `\_e\_l\_` with known letters `v`, `e`, `l` and excluded letters `s`, `a`, `t`, `o`, `i`, `r`, I will immediately stop listing invalid candidates. Instead, I will directly propose a single, valid English word that fits the constraints (e.g., "vile" is too short, so I must find a 5{-}letter word like "velum" if 'm' is valid, or "vella", or "velar" if 'a'/'r' weren't red, or "level" if positions allow). I will ensure the explanation is brief, explicitly states the logic for the *new* guess, and ends with a valid 5{-}letter word in the "guess:" tag, avoiding any repetition or rambling. Specifically, I will verify the word length and validity before outputting.} \\
\\ 
\\ 
\texttt{You are a language wizard who likes to guess words by using the given rules.} \\
\\ 
\texttt{Welcome to Wordle! You have six attempts to guess the target word, a valid English word of five lowercase letters (a{-}z). Please use the tags "explanation:" and "guess:" to provide a concise explanation for each guess.} \\
\\ 
\texttt{For instance, if your guess is "apple", your response should be} \\
\texttt{explanation: this is a common five{-}letter English word, and I am starting my guess with this word.} \\
\texttt{guess: apple} \\
\\ 
\texttt{After each guess, your answer will be validated, and you will receive feedback indicating which letters are correct (green), which letters are correct but in the wrong position (yellow), and which letters are incorrect (red). This feedback can be useful in determining which letters to include or exclude in your next guess.} \\
\\ 
\texttt{For example, the feedback for "apple" might be:} \\
\texttt{guess\_feedback: a<yellow> p<yellow> p<green> l<yellow> e<red>} \\
\\ 
\texttt{The explanation should contain details about how the guess\_feedback is used to arrive at a new guess.} \\
\\ 
\texttt{Do not generate any other text.} \\
\\ 
\texttt{Let's begin with your first guess.} \\
            }
        }
    }
     \\ \\

    \theutterance \stepcounter{utterance}  
    & \multicolumn{4}{p{0.6\linewidth}}{
        \cellcolor[rgb]{0.561,0.737,0.561}{
            \makecell[{{p{\linewidth}}}]{
                \texttt{\tiny{[P1$\rangle$GM]}}
                \texttt{explanation: I am starting with a common five{-}letter word containing frequent vowels and consonants to maximize initial information gain.} \\
\texttt{guess: slate} \\
            }
        }
    }
    & & \\ \\

    \theutterance \stepcounter{utterance}  
    & & & \multicolumn{2}{p{0.3\linewidth}}{
        \cellcolor[rgb]{0.9,0.9,0.9}{
            \makecell[{{p{\linewidth}}}]{
                \texttt{\tiny{[GM$|$GM]}}
                \texttt{attempts = 1} \\
\texttt{target\_word = level} \\
\texttt{guess = slate} \\
\texttt{guess\_feedback = s<red> l<yellow> a<red> t<red> e<yellow>} \\
            }
        }
    }
    & & \\ \\

    \theutterance \stepcounter{utterance}  
    & & & \multicolumn{4}{p{0.6\linewidth}}{
        \cellcolor[rgb]{0.9,0.9,0.9}{
            \makecell[{{p{\linewidth}}}]{
                \texttt{\tiny{[P1$\langle$GM]}}
                \texttt{guess\_feedback: s<red> l<yellow> a<red> t<red> e<yellow>} \\
\\ 
\texttt{Provide your response only in this format.} \\
\texttt{explanation: details} \\
\texttt{guess: word} \\
            }
        }
    }
     \\ \\

    \theutterance \stepcounter{utterance}  
    & \multicolumn{4}{p{0.6\linewidth}}{
        \cellcolor[rgb]{0.561,0.737,0.561}{
            \makecell[{{p{\linewidth}}}]{
                \texttt{\tiny{[P1$\rangle$GM]}}
                \texttt{explanation: The feedback shows 'l' and 'e' are in the word but misplaced, while 's', 'a', and 't' are excluded. I need a word with 'l' and 'e' in new positions, avoiding the red letters. "Clever" is too long, so I try "level" which fits the pattern and uses common letters.} \\
\texttt{guess: level} \\
            }
        }
    }
    & & \\ \\

    \theutterance \stepcounter{utterance}  
    & & & \multicolumn{2}{p{0.3\linewidth}}{
        \cellcolor[rgb]{0.9,0.9,0.9}{
            \makecell[{{p{\linewidth}}}]{
                \texttt{\tiny{[GM$|$GM]}}
                \texttt{attempts = 2} \\
\texttt{target\_word = level} \\
\texttt{guess = level} \\
\texttt{guess\_feedback = l<green> e<green> v<green> e<green> l<green>} \\
            }
        }
    }
    & & \\ \\

    \theutterance \stepcounter{utterance}  
    & & & \multicolumn{2}{p{0.3\linewidth}}{
        \cellcolor[rgb]{0.9,0.9,0.9}{
            \makecell[{{p{\linewidth}}}]{
                \texttt{\tiny{[GM$|$GM]}}
                \texttt{game\_result = WIN} \\
            }
        }
    }
    & & \\ \\

\end{supertabular}
}

\end{document}
